
\serie{Multiples et diviseurs}

\begin{exercice}[Vocabulaire]
Réponds aux questions suivantes en justifiant :
\begin{enumerate}
 \item 4 est-il un diviseur de 28 ?
 \item 32 est-il un multiple de 6 ?
 \item 4 divise-t-il 18 ?
 \item 35 est-il divisible par 5 ?
 \end{enumerate}
\end{exercice}

\begin{exercice}
Dans chaque cas, écris quatre phrases utilisant les nombres et l'un des mots suivants : diviseur, multiple, divisible, divise.
\begin{colenumerate}{3}     
 \item 70 et 210 ;
 \item 195 et 15 ;
 \item 192 et 48.
 \end{colenumerate}
\end{exercice}

\begin{exercice}[Critères de divisibilité]
Parmi les nombres : 12 ; 30 ; 27 ; 246 ; 325 ; 4\,238 et 6\,139, indique ceux qui sont divisibles :
\begin{colenumerate}{3}     
 \item par 2 ;
 \item par 3 ;
 \item par 5 ;
 \item par 9 ;
 \item par 10 ;
 \item par 25.
 \end{colenumerate}
\end{exercice}

\begin{exercice}
Parmi les nombres : 21 ; 12 ; 2 ; 619 ; 999 ; 416 ; 296 ; 540 ; 1\,785, quels sont les nombres divisibles par :
\begin{colenumerate}{3}     
 \item 2 ?
 \item 3 ?
 \item 5 ?
 \item 9 ?
 \item 10 ?
 \item 25 ?
 \end{colenumerate}
\end{exercice} 

\begin{exercice}
Parmi les nombres 15 ; 17 ; 58 ; 106 ; 54 ; 125 ; 105 ; 1\,577 ; 204, quels sont les nombres divisibles par :
\begin{colenumerate}{3}     
 \item 2 ?
 \item 3 ?
 \item 5 ?
 \item 9 ?
 \item 10 ?
 \item 25 ?
 \end{colenumerate}
\end{exercice} 

\begin{exercice} On s'intéresse aux nombres de trois chiffres de la forme $65 u$ où $u$ représente le chiffre des unités.

Quelles sont les valeurs possibles de $u$ pour obtenir :
\begin{enumerate}
 \item un multiple de 2 ?
 \item un nombre divisible par 3 ?
 \item un nombre divisible par 9 ?
 \end{enumerate}
\end{exercice} 

\begin{exercice}[Division et diviseurs]
\begin{enumerate}
 \item Effectue la division de 126 par 7 ;
 \item Déduis‑en deux diviseurs de 126.
 \end{enumerate}
\end{exercice} 

\begin{exercice}[Diviseurs]
\begin{enumerate}
 \item Écris deux nombres divisibles par 3 mais pas par 9 ;
 \item Écris deux multiples de 5 divisibles par 9 ;
 \item Écris le plus grand diviseur de 36 différent de 36.
 \end{enumerate}
\end{exercice} 

\begin{exercice}[Multiples]
\begin{enumerate}
 \item Trouve des multiples à la fois de 3 et de 5. Sont‑ils tous des multiples de 15 ?
 \item Trouve des multiples à la fois de 3 et de 6. Sont‑ils tous des multiples de 18 ?
 \end{enumerate}
\end{exercice} 

\begin{exercice}[Multiples (bis)]
\begin{enumerate}
 \item Écris trois multiples de 24 et quatre multiples de 18 ;
 \item Trouve le plus grand multiple de 12 inférieur à 75 et le plus grand multiple de 36 inférieur à 100 ;
 \item Cite un nombre multiple de 2 dont un  diviseur est 3.
 \end{enumerate}
\end{exercice} 

\begin{exercice}[Liste]
\begin{enumerate}
 \item Trouve tous les nombres divisibles par 5 compris entre 220 et 260 ;
 \item Parmi ces nombres, quels sont ceux qui sont divisibles par 3 ?
 \end{enumerate}
\end{exercice} 

\begin{exercice}[Énigme]
Trouve tous les nombres de trois chiffres divisibles à la fois par 3 et par 5 et dont le chiffre des centaines est 7.
\end{exercice} 

\begin{exercice}
Sur chacun des traits \_ , mettre un chiffre pour que les nombres composés soit divisibles par 3 et 5 :
\begin{enumerate}
 \item \_ $32$ \_ ;
 \item $4$ \_ $8$ \_ ;
 \item $2$ \_ \_ $5$.
 \end{enumerate}
\end{exercice}


\begin{exercice}[Décompositions]
\begin{enumerate}
 \item Décompose 18 sous la forme d'un produit de deux facteurs entiers différents de 1 ;
 \item Décompose 12 sous la forme d'un produit de trois facteurs entiers différents de 1 ;
 \item Peux‑tu décomposer 7 sous la forme d'un produit de deux facteurs entiers différents de 1 ? Et de 3 ?
 \end{enumerate}
\end{exercice} 

\begin{exercice}[Diviseurs communs à \ldots]
\begin{enumerate}
 \item Quels sont les diviseurs de 12 ? Cites-les tous ;
 \item Quels sont les diviseurs de 15 ? Cites-les tous ;
 \item Quels sont les diviseurs communs de 15 et de 12 ? Pourquoi ?
 \end{enumerate}
\end{exercice} 

\begin{exercice}
Écris la liste de tous les diviseurs de :
\begin{colenumerate}{4}
 \item 32 ;
 \item 67 ;
 \item 81 ;
 \item 144.
 \end{colenumerate}
\end{exercice} 

\begin{exercice}[Multiples communs à \ldots]
\begin{enumerate}
 \item Écris quelques multiples de 18. Peux‑tu les citer tous ?
 \item Écris quelques multiples de 15. Peux‑tu les citer tous ?
 \item Quels sont les multiples communs de 18 et de 15 ?
 \end{enumerate}
\end{exercice} 

\begin{exercice}[Encadrement]
\begin{enumerate}
 \item Encadre 55 puis 193 par des multiples consécutifs de 2 ;
 \item Encadre 56 puis 88 par des multiples consécutifs de 3 ;
 \item Encadre 125 puis 255 par des multiples consécutifs de 4.
 \end{enumerate}
\end{exercice} 

%%%%%%%%%%%%%%%%%%%%%%%%%%%%%%%%%%%%%%%%%%%%%%%%%%%%%%%%%%%%%%%%%%

\serie{Puissances}

\begin{exercice}[Vocabulaire]
Voici une liste de mots : base, exposant, puissance, facteurs, produit. Recopie chaque phrase en la complétant par le mot qui convient :
\begin{enumerate}
 \item $3^7$ se lit « 3 \ldots \ldots \ldots \ldots \ldots 7 » ;
 \item $5^4$ est le \ldots \ldots \ldots de quatre \ldots \ldots \ldots \ldots tous égaux à 5 ;
 \item 8 est l' \ldots \ldots \ldots de $6^8$ et 6 est la \ldots \ldots \ldots ;
 \item Le \ldots \ldots \ldots de six \ldots \ldots \ldots égaux s'écrit sous la forme d'une \ldots \ldots \ldots d'\ldots \ldots \ldots 6.
 \end{enumerate}
\end{exercice} 

\begin{exercice}[D'une écriture à l'autre]
\begin{enumerate}
 \item Écris en toutes lettres : $3^4$ ; $2^3$ et $7,1^9$ ;
 \item Écris en expressions mathématiques :
  \begin{colitemize}{2}
   \item huit puissance neuf ;
   \item quatre au cube ;
   \item trois puissance cinq ;
   \item sept au carré.  
   \end{colitemize}
 \end{enumerate}
\end{exercice}

\begin{exercice}[Notations puissance]
Recopie et complète chaque expression par le (ou les) exposant(s) manquant(s) :
\begin{enumerate}
 \item $4 \cdot 4 \cdot 4 \cdot 4 \cdot 4 \cdot 4 \cdot 4 \cdot 4 \cdot 4 = 4^{...}$ ;
 \item $0,1 \cdot 0,1 \cdot 0,1 = 0,1^{...}$ ;
 \item $2 \cdot 3 \cdot 2 \cdot 2 \cdot 3 = 2^{...} \cdot 3^{...}$ ;
 \item $2 \cdot 2 \cdot 3 \cdot 3 \cdot 3 \cdot 5 = 2^{...} \cdot 3^{...} \cdot 5^{...}$.
 \end{enumerate}
\end{exercice}

\begin{exercice}[Décomposition]
Décompose chaque expression en faisant apparaître des produits de nombres sans puissance :
\begin{colenumerate}{3}
 \item $2^4$ ;
 \item $7^2$ ;
 \item $0,1^5$ ;
 \item $1,2^2$ ;
 \item $3^2 \cdot 5^3$ ;
 \item $10^2 \cdot 6^1$.
 \end{colenumerate}
\end{exercice}

\begin{exercice}
Décompose puis donne l'écriture décimale sans l’aide de calculatrices :
\begin{colenumerate}{3}
 \item $2^4$ ;
 \item $7^2$ ;
 \item $0,1^3$ ;
 \item $1,2^2$ ;
 \item $1^5$ ;
 \item $0^4$ ;
 \item $17^1$ ;
 \item $12^0$ ;
 \item $10^3$.
 \end{colenumerate}
\end{exercice}


\begin{exercice}[Calculatrice]
Donne l'écriture décimale en calculant à la calculatrice :
\begin{colenumerate}{4}
 \item $2^{14}$ ;
 \item $17^7$ ;
 \item $8^{11}$ ;
 \item $1,2^6$.
 \end{colenumerate}
\end{exercice}

%%%%%%%%%%%%%%%%%%%%%%%%%%%%%%%%%%%%%%%%%%%%%%%%%%%%%%%%%%%%%%%%%%

\serie{Nombres premiers et décompositions}

\begin{exercice}[Produit de puissances]
Écris les nombres suivants sous la forme d'un produit :
\begin{enumerate}
 \item De puissances de 2 et de 5 :
 
$A = 2 \cdot 2 \cdot 5 \cdot 5 \cdot 5 \cdot 2 \cdot 2 \cdot 5 \cdot 5$ ;

$B = 25 \cdot 10 \cdot 5 \cdot 8$ ;

$C = 625 \cdot 512$ ;
 \item De puissances de 2, de 3 et de 7 :
 
$D = 2 \cdot 2 \cdot 2 \cdot 3 \cdot 7 \cdot 7$ ;

$E = 32 \cdot 21 \cdot 12$ ;

$F = 12 \cdot 21 \cdot 49$ ;

$G = 42$.
 \end{enumerate}
\end{exercice} 

\begin{exercice}[Drôles de dames \ldots]
Nefissa a décomposé 500 en un produit de 3 facteurs : $500 = 2 \cdot 25 \cdot 10$. Ivete affirme qu'elle peut décomposer $500$ avec plus de facteurs. Donne toutes les décompositions qu'Ivete a pu trouver.
\end{exercice} 

\begin{exercice}
Parmi les nombres suivants, quels sont ceux qui sont premiers ?
\begin{colitemize}{4}
 \item 25 ;
 \item 17 ;
 \item 36 ;
 \item 4 ;
 \item 99 ;
 \item 27 ;
 \item 56 ;
 \item 19 ;
 \item 12 ;
 \item 31 ;
 \item 88 ;
 \item 1.
 \end{colitemize}
\end{exercice} 

\begin{exercice}[Produit de facteurs premiers]
Donne la décomposition en facteurs premiers des nombres suivants :
\begin{colenumerate}{3}
 \item 96 ;
 \item 220 ;
 \item 245 ;
 \item 252 ;
 \item 360 ;
 \item 405 ;
 \item 420 ;
 \item 480 ;
 \item 891.
 \end{colenumerate}
\end{exercice} 


%%%%%%%%%%%%%%%%%%%%%%%%%%%%%%%%%%%%%%%%%%%%%%%%%%%%%%%%%%%%%%%%%%

\serie{Diviseurs communs, PGDC}

\begin{exercice}[Liste des diviseurs communs et PGDC]
Dans chaque cas, écris la liste des diviseurs communs aux deux nombres et entoure leur PGDC :
\begin{colenumerate}{3}
 \item 12 et 8 ;
 \item 8 et 10 ;
 \item 2 et 6 ;
 \item 6 et 9 ;
 \item 8 et 18 ;
 \item 12 et 20 ;
 \item 15 et 45 ;
 \item 27 et 18 ;
 \item 32 et 25 ;
% \item 56 et 84 ;
% \item 55 et 75 ;
% \item 124 et 1.
  \end{colenumerate}
\end{exercice} 

\begin{exercice}
Décompose en produits de facteurs premiers les nombres suivants. Donne le résultat en utilisant des puissances :
\begin{colenumerate}{3}
 \item 36 ;
 \item 2\,835 ;
 \item 17 ;
 \item 19 ;
 \item 660 ;
 \item 81.
  \end{colenumerate}
\end{exercice} 

\begin{exercice}
Décompose en produits de facteurs premiers les nombres suivants. Donne le résultat en utilisant des puissances :
\begin{colenumerate}{3}
 \item 69 ;
 \item 1\,125 ;
 \item 1225 ;
 \item 40 ;
 \item 37 ;
 \item 462.
 \end{colenumerate}
\end{exercice} 

\begin{exercice}
Trouve les nombres correspondant aux produits suivants :
\begin{colenumerate}{3}
 \item $2^3 \times 3^2 \times 5$ ;
 \item $4^2 \times 5^2$ ;
 \item $7^0 \times 11^2$.
 \end{colenumerate}
\end{exercice} 

\begin{exercice}
Décompose les nombres suivants en produits de facteurs premiers et trouve le PGDC :
\begin{colenumerate}{2}
 \item 36 et 48 ;
 \item 12 et 225.
 \end{colenumerate}
\end{exercice} 

\begin{exercice}
Décompose les nombres suivants en produits de facteurs premiers et trouve le PGDC :
\begin{colenumerate}{2}
 \item 384 et 1\,024 ;
 \item 105 et 180.
 \end{colenumerate}
\end{exercice} 

\begin{exercice}
Décompose les nombres suivants en produits de facteurs premiers et trouver le PGDC :
\begin{colenumerate}{2}
 \item 286 et 51 ;
 \item 84 et 198.
 \end{colenumerate}
\end{exercice} 

\begin{exercice}[Nombre de joueurs]
Dans un jeu, 180 jetons noirs et 120 jetons blancs doivent être tous répartis entre les joueurs. Tous les joueurs doivent avoir le même nombre de jetons noirs et le même nombre de jetons blancs.
\begin{enumerate}
 \item Peut-il y avoir vingt joueurs ? Neuf joueurs ?
 \item Combien peut-il y avoir de joueurs ? Donne toutes les possibilités.
 \end{enumerate}
\end{exercice} 

\begin{exercice}[Chez le fleuriste]

\vspace{1em}

\begin{minipage}[c]{0.66\linewidth}
Un fleuriste dispose de 30 marguerites et de 24 tulipes. Il veut utiliser toutes ses fleurs pour composer des bouquets tous identiques (chaque bouquet aura le même nombre de marguerites et le même nombre de tulipes que les autres).
 \end{minipage}\hfill%
 \begin{minipage}[c]{0.26\linewidth}
 \includegraphics[width=1.8cm]{fleurs}
  \end{minipage} \\[1em]
 Combien de bouquets peut-il réaliser ? Quelle est la composition de chaque bouquet ? Donne toutes les possibilités.
\end{exercice} 

\begin{exercice}[Le pâtissier]
Un pâtissier dispose de 75 pommes, 50 oranges et 100 poires. Afin de préparer des tartes, il désire répartir ces fruits en les utilisant tous et en obtenant le maximum de tartelettes identiques. Calcule le nombre de tartelettes et indique leur composition.
\end{exercice}

\begin{exercice}[Nombres croisés]
Recopie et complète la grille à l'aide des nombres que tu trouveras grâce aux définitions.

\begin{center}
\renewcommand*\tabularxcolumn[1]{>{\centering\arraybackslash}m{#1}}
\begin{tabularx}{.6\linewidth}{|X|X|X|X|X|} 
\cline{2-5}
\multicolumn{1}{c|}{} & \cellcolor {F2} A & \cellcolor {F2} B & \cellcolor {F2} C & \cellcolor {F2} D \\\hline
\cellcolor{H1} I & \cellcolor{G3} & \cellcolor{G3} & \cellcolor{G3} & \cellcolor{Noir} \\\hline
\cellcolor{H1} II & \cellcolor{G3} & \cellcolor{G3} & \cellcolor{G3} & \cellcolor{G3} \\\hline
\cellcolor{H1} III & \cellcolor{G3} & \cellcolor{Noir} & \cellcolor{G3} & \cellcolor{G3} \\\hline
\cellcolor{H1} IV & \cellcolor{Noir} & \cellcolor{G3} & \cellcolor{Noir} & \cellcolor{G3} \\\hline
\end{tabularx}
\end{center}

\textbf{Horizontalement} :

\textcolor{H1}{\textbf{I}} : PGDC $(125 ; 250)$ ;

\textcolor{H1}{\textbf{II}} : Ce nombre est un multiple de 9 ;

\textcolor{H1}{\textbf{III}} : Le chiffre des unités d'un nombre divisible par 10 -- Ce nombre est divisible par 5 ;

\textcolor{H1}{\textbf{IV}} : Le reste de la division euclidienne de 121 par 8 -- Le quotient dans celle de 245 par 112.

\textbf{Verticalement} :

\textcolor{F2}{\textbf{A}} : Le plus petit multiple de 24 à trois chiffres ;

\textcolor{F2}{\textbf{B}} : Le quotient de la division euclidienne de 274 par 10 -- Diviseur commun à tous les entiers ;

\textcolor{F2}{\textbf{C}} : Le chiffre des centaines est 5, celui des unités 4 et c'est le PGDC $(1\,542 ; 3\,598)$ ;

\textcolor{F2}{\textbf{D}} : 3 est un diviseur de ce nombre.
\end{exercice}

%%%%%%%%%%%%%%%%%%%%%%%%Mise en page
\newpage
%%%%%%%%%%%%%%%%%%%%%%%%%%%%%%%%%%%


\begin{exercice}[Carrelage]
Dans une salle de bain, on veut recouvrir le mur se trouvant au-dessus de la baignoire avec un nombre entier de carreaux de faïence de forme carrée dont le côté est un nombre entier de centimètres, le plus grand possible. Détermine la longueur, en centimètres, du côté d'un carreau de faïence sachant que le mur mesure 210 cm de hauteur et 135 cm de largeur. Combien faudra-t-il alors de carreaux ?
\end{exercice}

\begin{exercice}[Quel est l'âge d'Aloys ?]
Aloys adore les devinettes. Lorsqu'on lui demande son âge, cette année il répond :
\begin{itemize}
 \item L'an prochain , mon âge sera divisible par 2 ;
 \item Dans deux ans , mon âge sera divisible par 3 ;
 \item Dans trois ans, mon âge sera divisible par 4 ;
 \item Dans quatre ans, mon âge sera divisible par 5 ;
 \item Et j'ai moins de 97 ans.
 \end{itemize}
Quel est l'âge d'Aloys ?
\end{exercice}


%%%%%%%%%%%%%%%%%%%%%%%%%%%%%%%%%%%%%%%%%%%%%%%%%%%%%%%%%%%%%%%%%%

\serie{Multiples communs, PPMC}

\begin{exercice}[PPMC]
Dans chaque cas, donne le PPMC :
\begin{colenumerate}{3}
 \item 2 et 6 ;
 \item 12 et 8 ;
 \item 15 et 20 ;
 \item 20 et 30 ;
 \item 18 et 24 ;
 \item 48 et 36 ;
 \item 14 et 35 ;
 \item 18 et 20 ;
 \item 36 et 60.
 \end{colenumerate}
\end{exercice}

\begin{exercice}[Jeu de cartes]
On sait d'un jeu de cartes qu’il possède un nombre de cartes qui est un multiple de 4 et de 5.
\begin{enumerate}
 \item Peut-il y avoir cinquante cartes ? vingt cartes ?
 \item Combien le paquet peut-il avoir de cartes ? Donne toutes les possibilités inférieures à 100.
 \end{enumerate}
\end{exercice}

\begin{exercice}[Quelle heure était-il ?]
A Zermatt, cinq petits bus effectuent des boucles différentes à partir de la place de la gare. Les durées des trajets sont variables suivant les bus : \\[-1em]
\begin{itemize}
 \item Ligne A : trajet de 40 minutes ;
 \item Ligne B : trajet de 20 minutes ;
 \item Ligne C : trajet de 30 minutes ;
 \item Ligne D : trajet de 10 minutes ;
 \item Ligne E : trajet de 25 minutes.
 \end{itemize}
À 16h30 , un touriste japonais reconnait les 5 bus qu'il a photographiés le matin au même endroit. \\[0.5em]
Saurais-tu dire quelle heure il était alors ?
\end{exercice}
